% Research-agenda addition (attribution strand): what the licence meter counts.
% Inserted after the closing pointer of \S\ref{sec:one-settlement}.
\subsection{What the licence meter counts}\label{sec:he-licence-split}

The settlement above billed Carol for the $16$ files her skill touched. That
count stays in the example as a cap and as an illustration of a usage meter.
This subsection replaces it as the basis for settlement. A licence can pay
for four different observables. Each one can be moved by the licensor, and
each moves in a different way.

\begin{table}[H]
\centering
\small
\caption{Four billing rules for a skill licence. Each pays for an observable
that the licensor can change without changing the outcome.}
\label{tab:he-billing-rules}
\begin{tabularx}{\textwidth}{@{}l >{\raggedright\arraybackslash}X >{\raggedright\arraybackslash}X@{}}
\toprule
\textbf{Rule} & \textbf{What it pays for} & \textbf{How the licensor moves it} \\
\midrule
Per load & the skill's instructions entering the agent's context & a broader description, so the skill triggers on more tasks \\
Per file & reference files opened, or buyer files touched & a line telling the agent to read every reference first; more, smaller files \\
Shapley share & the average marginal contribution of each component & split one needed component into several needed pieces; duplicate a substitute \\
Leave-one-out & $v(N)-v(N\setminus\{i\})$ for each component $i$ & split a needed component; each piece is paid the whole marginal value \\
\bottomrule
\end{tabularx}
\end{table}

The first two rules read the run's trace. The trace is not independent of
the payee. A skill's text sits in the agent's context as instruction, so the
licensor partly writes the behaviour the meter records.

\begin{property}[Meter steering]\label{prop:he-meter-steering}
\setlength{\parskip}{4pt plus 1pt}
Let a licence pay $m(\tau)$, where $\tau$ is the trace of a run. The skill's
text changes the distribution of $\tau$. Suppose two versions of a skill have
the same distribution of settled outcomes but different expected meter
values. Then a licensor who maximizes payment prefers the version with the
higher meter, and the buyer pays the difference for no outcome.

A payment computed only from settled outcomes, under an assignment of skills
to tasks that the licensor cannot see or set, is unchanged by such edits.
\textbf{Status:} \Designed{} boundary on the licence contract. It is the payee's
case of Theorem~\ref{thm:folk}: here the party being paid also writes the
input to the signal.
\end{property}

The last two rules read outcomes, not traces, so they survive
Property~\ref{prop:he-meter-steering}. They fail differently. Fix a set $N$
of licensed components that may be made available to a task, each owned by
one licensor. Let $v(S)$ be the expected settled value when exactly the
components in $S$ are available, so $v(\emptyset)$ is the value with none
of them. A \emph{royalty rule} $\phi$ assigns each component a share of
$v(N)-v(\emptyset)$. The first two properties below are among the Shapley
fairness conditions used for data marketplaces; the third is the condition
under which their replication result holds~\pdcite{agarwal2019market}:
\begin{itemize}[leftmargin=1.4em,itemsep=1pt]
  \item \textbf{Balance:} the shares sum to $v(N)-v(\emptyset)$.
  \item \textbf{Symmetry:} two components that add the same value to every
  coalition receive equal shares.
  \item \textbf{Anonymity:} the rule reads $v$ and the component names only,
  not which licensor owns which component.
\end{itemize}
A rule is \textbf{split-proof} if no licensor can raise its total share by
re-packaging its own components into a game with the same outcomes: one
component cut into pieces that act only together, or one component copied.

\begin{pdclaim}{Theorem}{Licence-Split Impossibility}
\setlength{\parskip}{4pt plus 1pt}
\label{thm:he-licence-split}
No royalty rule that satisfies Balance, Symmetry and Anonymity is
split-proof on every game.
\end{pdclaim}
\begin{proof}
\setlength{\parskip}{4pt plus 1pt}
\emph{Complements.} Let a task succeed only when every component in $N$ is
available, with $v(N)=1$ and $v(S)=0$ otherwise. Every component is then
symmetric to every other, so Balance and Symmetry force a share of $1/n$
each. A licensor owning $k<n$ of them receives $k/n$. If it cuts one of its
components into two pieces that work only together, the game has $n+1$
symmetric components and the same outcomes. Its share becomes
$(k+1)/(n+1)$, and
\[
\frac{k+1}{n+1}-\frac{k}{n}=\frac{n-k}{n(n+1)}>0 .
\]
Anonymity is used here: the rule cannot see that two components share an
owner. With $n=2$ and $k=1$ the share rises from $1/2$ to $2/3$.

\emph{Substitutes.} Agarwal, Dahleh and Sarkar prove that when seller
identities are anonymized, balance and robustness to replication cannot hold
together~\pdcite{agarwal2019market}. In their example two sellers offer
identical features and receive $1/2$ each. One seller lists a copy. Three
identical components share the same total, $1/3$ each, so the copier now
holds $2/3$.
\end{proof}

\begin{pdclaim}{Theorem}{Leave-One-Out Royalties}
\setlength{\parskip}{4pt plus 1pt}
\label{thm:he-loo-royalty}
The leave-one-out rule $\phi_i=v(N)-v(N\setminus\{i\})$ is not balanced. A
licensor can multiply its leave-one-out royalty by cutting a needed
component into needed pieces, and a component with an identical substitute
is paid nothing.
\end{pdclaim}
\begin{proof}
\setlength{\parskip}{4pt plus 1pt}
In the complements game every component's leave-one-out value is $v(N)-0=1$,
so $n$ components are paid $n$ against a total of $1$. Cutting one component
into $j$ needed pieces gives each piece the value $1$, so its owner's royalty
rises from $1$ to $j$. If two components are identical, removing either
leaves the other, so each has leave-one-out value $0$.
\end{proof}

\begin{pdexample}{Carol Splits Her Skill}
\setlength{\parskip}{4pt plus 1pt}
A constructed task needs two licensed skills: Carol's lint skill $C$ and
Eve's type-stub skill $E$. The task settles green only when both are
available. The royalty pool is $12$\,\textsc{cr} per green task.

\emph{Shapley, before the split.} There are two orders of arrival. Each skill
completes the pair in one of them, so each share is $1/2$:
$12\times\tfrac12=6$\,\textsc{cr} each.

\emph{Shapley, after the split.} Carol publishes $C$ as two skills, $C_1$
(rules) and $C_2$ (fixer), which work only together. There are $3!=6$ orders.
Each of the three skills arrives last, and completes the set, in $2$ of
them, so each share is $2/6=1/3$. Carol holds $C_1$ and $C_2$:
$12\times\tfrac23=8$\,\textsc{cr}. Eve receives
$12\times\tfrac13=4$\,\textsc{cr}. The outcome is unchanged.

\emph{Leave-one-out.} Before the split each skill's marginal value is the
whole $12$\,\textsc{cr}, so the rule owes $24$\,\textsc{cr} against a pool
of $12$. After the split Carol is owed $12+12=24$ and Eve $12$, a total of
$36$.

\emph{Pooled by licensor.} The rule sees two players, Carol and Eve, before
and after the split. Each receives $6$\,\textsc{cr} both times.
\end{pdexample}

\paragraph{The escape: pool by licensor.}
The impossibility depends on Anonymity. Drop it: merge each licensor's
components into one player before applying any rule. A re-packaging inside
one licensor then leaves the game the rule sees unchanged, so no share can
move. The remaining attack is a false name: register a second licensor key
and move a component to it. Pooling therefore turns the split attack into a
Sybil attack on the licensor key (\S\ref{sec:keystone}). Whether it is
deterred is the question that \pdchapref{stp}{From Spawn to Person} prices
for reset identities%
\ifpdbook
\ (Theorem~\ref{stp:thm:whitewash-cost})%
\fi
. A fresh key enters under the newcomer ceilings, and the attack pays only
if its discounted gain exceeds the listing fee plus the earnings those
ceilings hold back. Computing that holdback for licence keys, whose earnings
are royalties rather than task income, is open. A licence with one licensor
needs no allocation at all: the split question arises only when several
licensors' components share one context.

\begin{definition}[Effect-gated royalty]\label{def:he-effect-royalty}
\setlength{\parskip}{4pt plus 1pt}
For each task that settles green and in which the harbor assigned the
licensor's components to be available, release
\[
  r \;=\; p\cdot\max\bigl(0,\;L\bigr),
\]
where $p$ is the agreed price per point of effect and $L$ is the lower
confidence bound, at a level fixed in the licence, on the licensor's pooled
availability effect. That effect is estimated from availability that the
harbor randomizes on comparable tasks, as \pdchapref{stp}{From Spawn to Person}
specifies%
\ifpdbook
\ (\S\ref{stp:sec:stp-skill-attribution})%
\fi
. The release stays in escrow through the dispute window and is clawed back
if a re-audit lowers $L$ or reverses the settlement. Loads, opens and
touched files are recorded as telemetry and may cap $r$. They never set it.
\textbf{Status:} \Designed{}; not implemented in the reference system.
\end{definition}

\noindent Take an illustrative price of $p=0.5$\,\textsc{cr} per point and a
lower bound of $L=12$ points. Carol's release is $0.5\times12=6$\,\textsc{cr}.
Kept as a cap, the per-file count allows $16\times0.5=8$, so the release is
$\min(8,6)=6$\,\textsc{cr}, and $10-6=4$\,\textsc{cr} returns to Bob. The
one-line check becomes $-456+400+50+6+0=0$. If $L\le0$, Carol receives
nothing however many files the skill opened.

\pdexercisepointer{\ref{ex:he-meter-steering-edit}}{\ref{ex:he-sybil-licensor}}{\pageref{ex:he-meter-steering-edit}}
